\Needspace{12\baselineskip}
\section{Fuerzas dependientes de la velocidad}

Una partícula de masa $m$ que se mueve en una dimensión bajo una fuerza que solo depende de la velocidad, $F=F(v)$, queda descrita por la segunda ley de Newton. Ejemplos son el arrastre lineal, $F=-bv$ (bajas velocidades), y el arrastre cuadrático, $F=-Bv^2$ (altas velocidades).

\begin{modelo}[movimiento con $F=F(v)$]
Hay que resolver las dos ecuaciones de primer orden
\begin{equation}
  m\,\dv{v}{t}=F(v),
  \label{eq:vel-newton}
\end{equation}
\begin{equation}
  \dv{x}{t}=v(t).
  \label{eq:vel-posicion}
\end{equation}
Equivalen a la ecuación de segundo orden $m\ddot x=F(\dot x)$, de modo que la solución general $x=\varphi(t,c_1,c_2)$ depende de dos constantes. Estas se determinan con dos condiciones iniciales: $x(t_0)=x_0=\varphi(t_0,c_1,c_2)$ y $\dot x(t_0)=v_0=\varphi'(t_0,c_1,c_2)$.
\end{modelo}

\begin{algoritmo}[movimiento con $F=F(v)$]
\begin{pasos}
  \item \emph{Velocidad.} Separar variables en \eqref{eq:vel-newton}: $\displaystyle\int\frac{\dd v}{F(v)}=\int\frac{\dd t}{m}$. Esto da $t=t(v)$ y, al invertir, $v=v(t)$ (constante $c_1$).
  \item \emph{Posición.} Integrar $\dd x=v(t)\,\dd t$, es decir, $x=x(t)$ (constante $c_2$).
  \item \emph{Condiciones iniciales.} Imponer $v(t_0)=v_0$ y $x(t_0)=x_0$.
  \item \emph{Variante (si interesa $x$ y no $t$).} Con la regla de la cadena, $\dot v=\dv{v}{x}\dv{x}{t}=v\dv{v}{x}$, y \eqref{eq:vel-newton} se escribe $m\,v\,\dd v=F(v)\,\dd x$, ecuación de variables separables que relaciona $x$ y $v$.
\end{pasos}
\end{algoritmo}

\begin{example}[arrastre lineal, $F(v)=-kv$]
Consideramos un movimiento horizontal sin otras fuerzas, con $v(0)=v_0$. Separando variables,
\[
  m\dv{v}{t}=-kv\;\Longrightarrow\;\frac{\dd v}{v}=-\frac{k}{m}\,\dd t .
\]
Si $v(t_f)=v_f<v_0$ (la partícula frena), integrando entre $0$ y $t_f$:
\[
  \int_{v_0}^{v_f}\frac{\dd v}{v}=-\frac km\int_0^{t_f}\dd t
  \;\Longrightarrow\;
  \ln\frac{v_f}{v_0}=-\frac km\,t_f
  \;\Longrightarrow\;
  t_f(v_f)=-\frac mk\ln\frac{v_f}{v_0}.
\]
Como $v_f/v_0<1$, $\ln(v_f/v_0)<0$ y $t_f>0$. Aplicando la exponencial a ambos lados de la ecuación logarítmica, $v_f/v_0=\ee^{-kt_f/m}$, es decir,
\begin{equation}
  v(t)=v_0\,\ee^{-kt/m}.
  \label{eq:vel-lineal}
\end{equation}
La velocidad decrece exponencialmente y tiende a cero, aunque no se anula en tiempo finito. La posición se obtiene integrando \eqref{eq:vel-posicion}:
\[
  x(t)=x_0+\frac{mv_0}{k}\bigl(1-\ee^{-kt/m}\bigr)\;\longrightarrow\;x_0+\frac{mv_0}{k}\quad(t\to\infty),
\]
de modo que la distancia total recorrida es finita, $mv_0/k$.\finejemplo
\end{example}

\begin{example}[arrastre cuadrático, $F(v)=-Bv^2$]
De $m\,\dd v/v^2=-B\,\dd t$ se obtiene $1/v=1/v_0+Bt/m$, es decir,
\[
  v(t)=\frac{v_0}{1+Bv_0t/m},\qquad x(t)=x_0+\frac mB\ln\!\Bigl(1+\frac{Bv_0t}{m}\Bigr).
\]
La velocidad decae como $1/t$ y el recorrido diverge logarítmicamente. Con la variante del algoritmo: $m\,v\,\dd v/\dd x=-Bv^2$ da $\dd v/v=-(B/m)\,\dd x$ y $v=v_0\,\ee^{-B(x-x_0)/m}$.\finejemplo
\end{example}

El caso $F=-k/v^2$, en el que la partícula se detiene en tiempo finito, se resuelve con las dos variantes del algoritmo en el problema~3 de la \cref{sec:seminarios}.

\begin{figure}[H]
  \centering
  \begin{tikzpicture}
    \begin{axis}[informe, xmin=0, xmax=3, ymin=0, ymax=1.05,
      xlabel={$t$ (con $m=k=v_0=1$)}, ylabel={$v/v_0$}, legend style={at={(0.97,0.95)}, anchor=north east}]
      \addplot[serie1, domain=0:3, samples=100] {exp(-x)};
      \addlegendentry{$F=-kv$}
      \addplot[serie2, domain=0:3, samples=100] {1/(1+x)};
      \addlegendentry{$F=-Bv^2$}
      \addplot[serie4, domain=0:0.3333, samples=100] {(1-3*x)^(1/3)};
      \addlegendentry{$F=-k/v^2$}
    \end{axis}
  \end{tikzpicture}
  \caption{Frenado con distintas fuerzas dependientes de la velocidad. Con $F=-k/v^2$ la partícula se detiene en el instante finito $t=mv_0^3/(3k)$.}
  \label{fig:velocidad}
\end{figure}

\newpage
